% Generated draft: requires mathematical and source review.
\documentclass[11pt]{article}
\usepackage[margin=25mm]{geometry}
\usepackage{fontspec}
\setmainfont{Latin Modern Roman}
\usepackage{amsmath,amssymb,mathtools,unicode-math}
\setmathfont{Latin Modern Math}
\usepackage{enumitem,needspace,xurl,hyperref}
\setlength{\parindent}{0pt}
\setlength{\parskip}{5pt}
\setlength{\emergencystretch}{3em}
\newcommand{\N}{\mathbb{N}}
\newcommand{\Z}{\mathbb{Z}}
\newcommand{\Q}{\mathbb{Q}}
\newcommand{\R}{\mathbb{R}}
\newcommand{\C}{\mathbb{C}}
\newcommand{\F}{\mathbb{F}}
\newcommand{\A}{\mathbb{A}}
\newcommand{\PP}{\mathbb{P}}
\newcommand{\E}{\mathbb{E}}
\newcommand{\eps}{\varepsilon}
\newcommand{\dd}{\,\mathrm{d}}
\begin{document}
\section*{UnsolvedMath ID: 30001679}
\textbf{Source problem code:} catalogue entry
\subsection{Definitions and mathematical statement}
Work in set theory with the axiom of choice. Write \(\omega_1\) for the first uncountable ordinal and \(\aleph_1=|\omega_1|\). For the formulation below, use the usual convention that an ideal \(I\) on \(\omega_1\) is proper, contains every bounded subset of \(\omega_1\), and is closed under taking subsets and countable unions. The supplied statement does not explicitly specify these ideal conventions, so they require confirmation against the original problem discussion.

For \(A,B\subseteq\omega_1\), put \(A\sim_I B\) if \(A\mathbin{\triangle}B\in I\). Forcing with the nonzero elements of \(\mathcal P(\omega_1)/I\), ordered by inclusion modulo \(I\), produces a generic ultrafilter \(U_G\) on the ground-model subsets of \(\omega_1\). The ideal \(I\) is \emph{precipitous} if this forcing forces the generic ultrapower \(\operatorname{Ult}(V,U_G)\), formed using ground-model functions, to be well-founded. The ideal \(I\) is \emph{normal} if, for every sequence \(\langle A_\xi:\xi<\omega_1\rangle\) of members of \(I\), its diagonal union
\[
\nabla_{\xi<\omega_1}A_\xi
=\{\alpha<\omega_1:\exists\xi<\alpha\ (\alpha\in A_\xi)\}
\]
belongs to \(I\).

Let \(\mathsf T\) denote the assertion
\[
\begin{gathered}
\text{There exists a precipitous ideal on }\omega_1,\\
\text{and there exists no ideal on }\omega_1
\text{ that is both normal and precipitous.}
\end{gathered}
\]
The second clause quantifies over all ideals on \(\omega_1\); it does not merely require the ideal witnessing the first clause to be nonnormal. The generalized continuum hypothesis, abbreviated GCH, is
\[
\forall\lambda\;\bigl(\lambda\text{ is an infinite cardinal}
\Longrightarrow 2^\lambda=\lambda^+\bigr).
\]
Here \(\kappa^{++}\) denotes the second successor cardinal of \(\kappa\).

The supplied problem asks:
\begin{enumerate}
\item Is the assumption \(o(\kappa)=\kappa^{++}\) needed to obtain a model satisfying \(\mathsf T\)? The notation \(o(\kappa)\) is conventionally interpreted in this context as the Mitchell order of a measurable cardinal \(\kappa\), but the supplied record does not define it or specify the underlying construction, other hypotheses, or the sense in which the assumption is claimed to be necessary. Thus this question is retained as a question about the required hypothesis, rather than converted into an unsupported implication or equiconsistency assertion.
\item Is there a model of \(\mathrm{ZFC}+\mathrm{GCH}+\mathsf T\)? The source does not specify the hypotheses from which such a model is to be constructed or its intended relative consistency strength.
\end{enumerate}
The supplied dated assessment labels the problem partially solved and reports a later construction satisfying \(\mathsf T\), while describing the assumption question and unspecified higher-cardinal variants as remaining open. These are supplied literature qualifications, not independently verified mathematical conclusions. No higher-cardinal variant is formulated in the supplied problem text.
\subsection{Short English statement}
The problem asks whether the hypothesis \(o(\kappa)=\kappa^{++}\) is needed to construct a model with a precipitous ideal on \(\aleph_1\) and no normal precipitous ideal on that cardinal, and whether such a model can satisfy GCH. The intended necessity claim and its background hypotheses are unspecified. The supplied partially solved label is dated curation requiring independent review.
\subsection{Sources}
\begin{itemize}
\item Supplied local UnsolvedMath registry record, permanent ID \(30001679\), in the Oberwolfach Reports Open Problems collection; exact source problem code: \texttt{catalogue entry}. The record gives the proposed year as 2011 and does not identify a proposer.
\item Supplied report citation: \emph{Set Theory} (2011), Oberwolfach Reports, pp.~85--140, DOI \texttt{10.4171/owr/2011/02}. Source URL and literature-source entry labelled ``Oberwolfach problem statement'': \url{https://doi.org/10.4171/owr/2011/02}.
\item Supplied literature-source entry labelled ``Gitik, precipitous ideals without normality'': \url{https://doi.org/10.1142/S0219061312500080}. The record attributes to later forcing work a construction of a precipitous ideal on \(\aleph_1\) without normal precipitous ideals; the theorem and its hypotheses have not been independently checked here.
\item Literature assessment checked at \texttt{2026-08-21}: the supplied status is \texttt{partially\_solved}. The supplied assessment reports that a central existence case was resolved while the listed assumption and higher-cardinal variants remain open. This is supplied dated curation, not independently verified literature assessment; the background explicitly requests independent verification of time-sensitive status and cited claims.
\item The supplied statement status is \texttt{corrected\_verified}; its accompanying verification says only that TeX notation was enclosed in mathematical delimiters without changing the mathematical content. This does not establish mathematical verification of the formulation or literature assessment.
\end{itemize}
\end{document}
